\documentclass[11pt,a4paper]{article}
\usepackage[utf8]{inputenc}
\newcommand{\undertitle}{Technical Appendix}
\newcommand{\headeright}{Technical Appendix}
\usepackage{emergentwisdom-preprint}
\usepackage{mathtools}
\usepackage{amssymb}
\usepackage{booktabs,tabularx,array,float}
\usepackage{enumitem}
\usepackage{xcolor}
\usepackage{microtype}
\usepackage{hyperref}
\newcolumntype{Y}{>{\raggedright\arraybackslash}X}
% Canonical compact interface shared with the Life Simulation snapshot.
% Keep this file theory-only: no implementation or experiment claims.
\newcommand{\MMCoreSummary}{%
Six forms carry the world model: Concept, Thing, Event, Binding, Cut, and
Realization. Bindings give Things roles in Events, while typed links connect
records without turning every connection into the same relation. A Cut answers
one question by dividing one stated unit among exclusive alternatives and a
remainder; those shares total one. A Realization links a Concept to an
addressable fragment without a weight, and any graded fit is expressed by a
separate Cut on an assessment Event. The six forms keep identity, occurrence,
participation, decomposition, and grounding apart without claiming to exhaust
metaphysics. A construction profile may add lifecycle, admissibility, and
recomposition rules while retaining the same forms. Understanding Nodes and
Document Nodes share their address space with world records but keep distinct
authority and validation, and perspective follows declared authority ancestry rather than a
duplicated attribution field.%
}

\newcommand{\MMCoreSchema}{%
\begin{center}
\footnotesize\ttfamily
\begin{tabular}{@{}l@{}}
Concept \{ id, names, typed relations, constraints,\\
\qquad grounding records, provenance \}\\
Thing \{ id, kind, identity boundary, continuity criterion,\\
\qquad typed links, provenance \}\\
Event \{ id, kind, interval, typed world data, optional description,\\
\qquad direction, typed links, provenance \}\\
Binding \{ id, event, role, thing, interval, provenance \}\\
Cut \{ id, parent Event, question, profile, divided unit,\\
\qquad keyed weighted answers, keyed remainder, optional conditioning address,\\
\qquad constraints, provenance \}\\
Realization \{ id, purpose, context parent, target, concept,\\
\qquad grounding version, optional alignment, evidence, provenance \}
\end{tabular}
\end{center}
\noindent\footnotesize Answer and remainder addresses pair the Cut revision
with a local key; they are not additional record forms.\normalsize%
}

\newcommand{\MMProfileSummary}{%
In the selected profile, each Thing has one lifecycle Event spanning its
modeled existence. Active joint Events and their Bindings determine
encapsulation, so no second containment truth is stored. Every Cut declares the
unit it partitions, gives weights to exclusive answers and a remainder, and
totals one; sequential weights come from covered duration. Understanding and
Document Nodes can point into the same graph as the six world forms without
acquiring world authority. An Event description is optional text on the Event,
not a Document Node.%
}

\newcommand{\MMProfileSchema}{%
\MMCoreSchema
\begin{center}
\footnotesize\ttfamily
Profile Cut conventions \{ allocation, sequential, direction;\\
\qquad sequential shares derived from duration \}
\end{center}%
}

\hypersetup{
  hidelinks,
  pdftitle={The Meaning Model: Minimized Grammar},
  pdfauthor={Henrik Westerberg},
  pdfsubject={Technical appendix specifying record forms, profile axioms, perspective projections, surface records, and interfaces},
  pdfkeywords={meaning model, grammar, specification, cut, remainder, lifecycle, containment, realization}
}
\renewcommand{\shorttitle}{Meaning Model Grammar}
\setlist[itemize]{leftmargin=1.4em}
\setlist[enumerate]{leftmargin=1.6em}
\setlist[description]{leftmargin=2.6em,labelindent=0pt,style=sameline,font=\normalfont\bfseries}
\newcommand{\Cset}{\mathcal C}
\newcommand{\Xset}{\mathcal X}
\newcommand{\Eset}{\mathcal E}
\newcommand{\Bset}{\mathcal B}
\newcommand{\Kset}{\mathcal K}
\newcommand{\Rset}{\mathcal R}
\newcommand{\Lset}{\mathcal L}
\newcommand{\lc}{\ell}
\newcommand{\tv}{\operatorname{TV}}
\title{The Meaning Model:\\Minimized Grammar}
\author{Henrik Westerberg\\Emergent Wisdom\\
  \texttt{henrik.westerberg@emergentwisdom.org}}
\ewpublicationdate{\today}
\begin{document}
\maketitle

\begin{abstract}
This appendix is a compact restatement of the definitions and rules explained
in \emph{The Meaning Model: Constructing Worlds and Stories at Progressive
Resolution}. It adds no requirements: the paper is self-contained, and this
reference groups the same account into record forms, obligations, selected-profile
axioms, perspective queries, surface records, and operations. Book-specific
choices remain non-normative, and profile rules bind only constructions that
select them. Both presentations specify the target contract, not the current
Rust host schema. Cut weights are the grammar's only semantic numbers;
externally grounded measurements remain typed Event data rather than another
weighted relation.
\end{abstract}

\keywords{Meaning Model, Grammar, Specification, Cut, Remainder, Lifecycle,
Containment, Realization}

\paragraph{Reading map.}
The paper's \emph{Grammar and Progressive Construction} section explains the
records, context rules, lifecycle profile, and refinement contract collected
here. \emph{Concepts and Examples} covers grounding and its operators;
\emph{World, Understanding, and Story} covers projections, access, surfaces,
and transactions. The Book profile and Appendix A explain the non-normative
person vocabulary and constants. A substantive rule change must be reflected
in both presentations; this reference is not a separate source of theory.

\section{Conventions}

\begin{description}[style=nextline]
  \item[Model.] A Meaning Model at declared resolution $\rho$ is a finite
    tuple $\mathcal M_\rho=(\Cset,\Xset,\Eset,\Bset,\Kset,\Rset,\Lset)$ of
    Concepts, Things, Events, Bindings, Cuts, Realizations, and typed links.
  \item[Shared address space.] A complete construction is the typed graph
    $\mathcal G_\rho=(\mathcal M_\rho,\mathcal U,\mathcal D,
    \mathcal L_{\times})$, where $\mathcal U$ contains Understanding Nodes,
    $\mathcal D$ contains Document Nodes, and $\mathcal L_{\times}$ contains
    links within and across the world, understanding, and document surfaces. All three surfaces use
    one identifier space. Membership in the same graph permits exact reference;
    it does not give an Understanding Node, Document Node, and Event the same authority, clock,
    composition rule, or validation contract.
  \item[Identity.] Every record reference identifies a logical identity and an
    exact revision; $\mathrm{id}$ below abbreviates the immutable revision
    address. Under the selected
    construction profile, a change creates a new revision address and a
    \texttt{supersedes} link (A11); the logical identity of a Thing or Concept
    does not thereby change. A Concept identifier is
    stable across grounding versions. A profile may identify a grounding
    version $g$ by the content hash of its declared bundle serialized in
    canonical order, as the Book does with Sema; a Realization's $g$ names the
    exact bundle it used. Content addressing is an identity convention, not a
    seventh record form or a universal requirement of the core grammar.
  \item[Provenance.] Every record carries
    $p=(\text{author or generator identifier},\ \text{commit},\ \text{construction time},\
    \text{evidence references})$.
  \item[Clocks.] Story time $T_s$ (intervals of Events), construction time
    $T_c$ (commits), discourse order $T_d$ (Document Node positions). A record
    states which clock each of its times uses.
  \item[Intervals.] A resolved interval is $I=[t_0,t_1)\subset T_s$.
    An ongoing end or unknown bound is explicitly tagged; duration
    $|I|=t_1-t_0$ is computed only for two finite resolved bounds.
    A moment Event has a tagged instant $t$ with support $\{t\}$, not a positive duration.
    Unresolved interval checks return \textsf{untestable}, not \textsf{pass}.
    Imagined and counterfactual timelines name their own temporal frame;
    comparisons and temporal containment require the same frame.
  \item[Roots (selected profile).] The Universe $U\in\Xset$ is a Thing and
    History $H=\lc(U)\in\Eset$ is its lifecycle Event. Every accepted-world
    lifecycle and Event is contained by History. Time is expressed through
    History and the intervals of its descendants, not by a separate Thing. A
    candidate is not part of History because it carries a date; it enters
    only through an accepted commit (A11).
\end{description}

\section{Record forms}

\MMCoreSchema

\begin{description}[style=nextline]
  \item[Concept] $c=(\mathrm{id},N,\mathrm{Rel},\mathrm{Constr},G,p)$.
    $N$ names; $\mathrm{Rel}\subseteq \mathrm{Kind}\times\Cset$ typed
    relations (specialization, opposition, expression, and other declared
    kinds); $\mathrm{Constr}$ constraints; $G$ grounding records, each of
    depth relational, representative prototype, exemplar, or canonical model,
    each versioned.
  \item[Thing] $x=(\mathrm{id},\mathrm{kind},\partial x,\kappa_x,\Lset_x,p)$.
    $\partial x$ identity boundary; $\kappa_x$ continuity criterion.
  \item[Event] $e=(\mathrm{id},\mathrm{kind},I_e,S_e,d_e,D_e,\Lset_e,p)$.
    State $S_e$ is a finite map from fields to categorical
    conditions or typed world data. A numeric datum is admissible only when it
    has an independently interpretable unit or measurement protocol, such as a
    date, duration, count, currency amount, distance, or physical measurement;
    measurement uncertainty may accompany it. The optional text $d_e$ is a
    readable description of this exact Event record and version;
    it is not a Document Node and does not change the structured record.
    A candidate description acquires accepted-world authority only if its
    candidate is committed.
    Interpretive judgments such as
    trust, intelligence, importance, and concept fit are Events and Cuts, not
    numerically anchored entries in $S_e$. Direction $D_e$ is a typed
    continuation specification (law, constraint set, predictor, or a reference
    to a direction Cut) or $\bot$; every numerical forecast allocation is a
    direction Cut.
  \item[Binding] $b=(\mathrm{id},e,r,x,I_b,p)$ with $I_b\subseteq I_e$; $r$ a typed
    role.
  \item[Cut] $K=(\mathrm{id},P,q,\pi,u_K,
    \{(k_i,a_i,w_i)\}_{i=1}^{n},(k_\bot,w_\bot),h,\mathrm{Constr},p)$. Parent
    $P\in\Eset$ is an Event; $q$ the question;
    $\pi\in\{\textsf{allocation},
    \textsf{sequential},\textsf{direction}\}$; children $a_i$ as fixed by
    $\pi$ (Axiom A6); $u_K$ the one divided unit. Distinct local keys
    $k_i,k_\bot$ make each answer and the remainder addressable as
    $(\mathrm{id},k)$, including zero-weight slots. These are addresses within
    a Cut revision, not additional world records. Optional $h$ names the
    answer whose share this Cut conditions on; $P$ remains its Event parent.
    A judgment about another
    record or fragment uses an assessment Event contained by the process whose
    perspective it expresses as $P$ and an \texttt{about} link to its target.
  \item[Realization] $r=(\mathrm{id},\mathrm{purpose},P,\tau,c,g,\sigma,
    \mathrm{ev},p)$. $\mathrm{purpose}\in\{\textsf{define},\textsf{describe}\}$;
    $P$ is the Event or surface node that supplies its structural context;
    $\tau$ an addressable subgraph of $\mathcal M_\rho$; $g$ a grounding
    version of $c$; $\sigma$ optional role alignment. The
    correspondence is unweighted. A graded judgment of fit is a separate Cut
    on an assessment Event under A10.
\end{description}

\paragraph{Typed links.}
\begin{center}
\footnotesize
\begin{tabularx}{\textwidth}{@{}l l Y@{}}
\toprule
Link & Domain $\times$ codomain & Meaning \\
\midrule
\texttt{contains} & Event $\times$ Event & temporal containment of occurrences in one frame \\
\texttt{occurs-within} & Event $\times$ Event & episode located inside a lifecycle or period \\
\texttt{context-parent} & record $\times$ context root or record & unweighted authority ancestry; not a claim of temporal inclusion \\
\texttt{updates} & Event $\times$ Event field or condition & unweighted assertion that the first Event changes the named datum or condition of the second \\
\texttt{about} & record $\times$ record & reference without participation or state change \\
\texttt{supersedes} & record $\times$ record & predecessor preserved; lineage changes; carries scope and clock \\
\texttt{realizes-forecast} & Event $\times$ Cut slot & identifies the selected direction answer, including the remainder \\
\bottomrule
\end{tabularx}
\end{center}
A multi-Thing relation is written as a joint Event, with a Binding for every
role; the grammar does not author direct edges between Things. Outside Event
mereology, records use unweighted context-parent
links. Context roots have labels such as \textsf{accepted-world}, \textsf{inner},
\textsf{understanding}, \textsf{document}, and \textsf{candidate}; authority
inheritance stops at the nearest such root. A model-context Event
may additionally declare \textsf{remembered}, \textsf{imagined}, or
\textsf{counterfactual} mode and a source reference. Its represented Events
need not occur during the mental act that hosts this context. Thus an inner
Event can be connected to History through a person's life without becoming a
public-world claim.

Only a Cut creates weighted parenthood: it partitions one stated unit among
exclusive sibling answers and a remainder. Event containment and ordinary
typed links remain unweighted, so linked Events may overlap or coexist. Event
payloads with external units are data, not a third kind of parent relation.

\section{Core obligations and selected-profile axioms}

A3--A5, A10, and A12 are core obligations of every Meaning Model. A1, A2,
A6--A9, A11, and A13 are axioms of the selected construction profile. The
labels below keep that boundary visible: using the six records alone does not
commit another construction to this lifecycle, time, or transaction policy.

\begin{description}
  \item[A1 Profile---Lifecycle.] For every $x\in\Xset$ there is exactly one
    $\lc(x)\in\Eset$ of kind \textsf{lifecycle} with Binding
    $(\lc(x),\textsf{subject},x,I_{\lc(x)})$. $I_{\lc(x)}$ is the modeled
    lifetime in the selected branch's active view. Superseded lifecycle
    revisions remain addressable but are not additional current lives.
    All state of $x$ is carried by $\lc(x)$ or by Events bound to $x$.
    The declared continuity criterion $\kappa_x$ decides whether a change
    preserves this Thing or creates a successor; split and merge Events bind
    the distinct identities they connect.
  \item[A2 Profile---Agency boundary.] For every Binding $(e,r,x,I_b)$ with $r$ an
    ordinary participant role, $I_b\subseteq I_{\lc(x)}$.
    Explicit \textsf{created} and \textsf{terminated} boundary roles may bind
    the appropriate lifecycle endpoint; they do not imply activity beyond
    it. A moment participant must lie inside its life interval. Other
    references outside that interval use \texttt{about}.
  \item[A3 Core---Perspective by containment.] Every Cut has an Event parent
    $P$. A perspective-dependent belief, appraisal, estimate, or judgment is a
    distinct Event beneath a declared perspective root. The Event's Cuts
    inherit that context from $P$. Only edges explicitly declared
    \textsf{authority-parent} participate in this inheritance: context-parent
    links and any containment links assigned that role by the profile.
    Every path must resolve to the same nearest context root, or the record
    is rejected as ambiguous. Occurrence, reference, evidence, and grounding
    links alone confer neither authority nor access. Profiles define
    admissible roots and paths;
    the selected lifecycle profile requires a person's perspective process to
    descend through that person's lifecycle, while a constructor or evaluator
    process descends from its named understanding-process root. A joint or
    public Event remains in shared History and identifies people or other
    Things through Bindings. No independently authored perspective field may
    contradict these structural paths.
  \item[A4 Core---Local simplex.] For every Cut,
    \begin{equation}
      w_i\ge 0,\qquad w_\bot\ge 0,\qquad \sum_{i=1}^{n} w_i+w_\bot=1,
    \end{equation}
    the children are mutually exclusive as shares of the declared unit
    $u_K$, and a record appearing under two Cuts has one weight in each.
    There is no global mass. Exclusivity concerns the allocated portions,
    not whether the referenced Events can coexist. An attribution over
    interacting causes must declare how shared responsibility is assigned,
    for example to an explicit joint answer; normalization alone cannot
    resolve double counting. Conditioning addresses must resolve, form no
    cycle, and identify the precise share whose conditional unit is divided.
  \item[A5 Core---Event processes and world data.] When a longer Event-process is
    opened into shorter happenings without a divided unit, the shorter Events
    are connected by \texttt{contains}, \texttt{occurs-within}, or
    \texttt{updates}; \texttt{about} is a reference, not a part relation.
    These links carry no weights, need not
    be exclusive, and assert no numerical recomposition. A numerical allocation
    over the same Events requires a separate Cut naming its question and unit.
    Events may also carry categorical facts and externally interpretable
    measurements. Such data record what happened; they do not assign semantic
    degree, importance, or psychological intensity.
  \item[A6 Profile---Named Cut families.]
    \begin{description}[leftmargin=1.2em]
      \item[allocation] Children reference existing records or items of a
        frozen vocabulary $V_q$. Named questions: \emph{facet}
        ($a_i\in V_q$, closed), \emph{contribution} ($a_i$ among Things
        bound to $P$), \emph{cause} ($a_i$ among existing wants, beliefs,
        actions, circumstances, or Events; $u_K$ is explanatory
        responsibility in the parent assessment's perspective).
      \item[sequential] Children $e_1,\dots,e_n\in\Eset$ with
        $I_{e_i}\subseteq I_P$ pairwise disjoint and ordered;
        \begin{equation}
          w_i=\frac{|I_{e_i}|}{|I_P|},\qquad w_\bot=1-\sum_i w_i .
        \end{equation}
        Shares are derived, never elicited. The parent interval must have
        finite positive duration. An open lifecycle is partitioned only over
        an explicitly bounded observation Event; an infinite or unresolved
        denominator is never used.
      \item[direction] Children are candidate continuations of $P$ at a
        declared horizon $\tau$; they are not parts of $P$. Weights are
        declared either a proposal policy or a forecast probability
        distribution before observing the outcome. A proposal policy guides
        construction; it does not assert world likelihood. Forecasts can be
        scored before they are well calibrated; calibration is an empirical
        property, not permission to call them probabilities. Alternatives
        must be exclusive at the stated horizon. At most one answer is
        realized, as a new Event with a \texttt{realizes-forecast} link to its slot;
        unrealized siblings are retained. Drawing the remainder requests an
        unnamed continuation, not silent renormalization over named answers.
    \end{description}
  \item[A7 Profile---Encapsulation and inherited presence.] Write
    $B(e,r,x,t)$ when a Binding assigns $x$ role $r$ in $e$ and $t$ lies in
    both the Event and Binding intervals. For $x\in\Xset$ and $t\in T_s$,
    \begin{equation}
    \begin{aligned}
      \operatorname{Enc}(x,t)=\bigl\{\,y:\ &\exists e\in\Eset\ [\
      (\mathrm{kind}(e)=\textsf{physical-part}\land B(e,\textsf{whole},x,t)\land B(e,\textsf{part},y,t))\\
      &\lor(\mathrm{kind}(e)=\textsf{membership}\land B(e,\textsf{group},x,t)\land B(e,\textsf{member},y,t))\\
      &\lor(\mathrm{kind}(e)=\textsf{presence}\land B(e,\textsf{place},x,t)\land B(e,\textsf{occupant},y,t))]\,\bigr\}.
    \end{aligned}
    \end{equation}
    $\operatorname{Enc}$ is unweighted. Rosters and co-presence are defined
    from the same Bindings. No containment among Things is authored directly.
    A location-propagating physical-part relation may derive the part's coarse
    presence from the whole until a detach Event ends that relation. When
    extents are modeled, the child's extent lies inside the whole's at the
    profile's declared tolerance. Social, institutional, and documentary
    membership never propagates location.
  \item[A8 Profile---Periods and episodes.] The periods of $x$ are the children of one
    authoritative sequential Cut of a bounded lifecycle interval, using an
    observation Event when the whole life is not resolved. An episode is
    an Event with an \texttt{occurs-within} link to $\lc(x)$ or to a period;
    episodes may overlap. Phases are children of sequential Cuts of episodes.
  \item[A9 Profile---Mixture law.] Let $K$ be a sequential Cut of a finite,
    positive-duration $P$ with children
    $e_i$, and let $F_i$ be facet Cuts on the $e_i$ with the same question,
    unit, vocabulary, remainder meaning, and governing perspective.
    If the children fully cover $P$, the facet vector is
    \begin{equation}
      f_P=\sum_i \lambda_i f_i,\qquad
      \lambda_i=\frac{|I_{e_i}|}{|I_P|}.
    \end{equation}
    If uncovered duration has share $\lambda_\bot>0$, the complete vector is
    defined only when an explicit compatible facet vector $f_\bot$ describes
    that uncovered interval, in which case
    $f_P=\sum_i\lambda_i f_i+\lambda_\bot f_\bot$. Otherwise the derived view
    reports the known partial contribution $\sum_i\lambda_i f_i$, coverage
    $1-\lambda_\bot$, and temporal remainder; it does not invent or claim a
    complete parent vector.
    When $F_P$ is already committed, even partial detail must leave a feasible
    residual:
    \begin{equation}
      r=f_P-\sum_i\lambda_i f_i\ge 0\quad\text{componentwise},\qquad
      \sum_v r_v=\lambda_\bot.
    \end{equation}
    If $\lambda_\bot>0$, $r/\lambda_\bot$ is the required residual
    composition, not evidence about the unobserved interval. For full
    coverage $r=0$ recovers the mixture equation. A declared numerical
    tolerance may bound these tests; it cannot authorize hidden
    renormalization. Infeasible detail is rejected or submitted as an
    explicit revision of $F_P$. Tolerances are validator settings, not world
    facts. The general conservative-opening contract is given below the
    interface table; A9 is its temporal specialization.
  \item[A10 Core---Realization.] A Realization is an unweighted
    correspondence between an addressable fragment and a Concept at grounding
    version $g$. \textsf{define} associates $c$ with an idealized fragment
    written in the record forms of Section~2; \textsf{describe} requires the
    grounding version used. The optional $\sigma$ records the role alignment.
    When graded fit is required, create an assessment Event beneath the
    relevant declared perspective process, link it \texttt{about} the target, and
    attach a Cut under the question
    ``how does one unit of this assessment of fit to $c$ divide?'' Its
    mutually exclusive answers are \textsf{matches}, \textsf{does-not-match},
    and remainder:
    \begin{equation}
      w_{\mathrm{match}}+w_{\mathrm{nonmatch}}+w_\bot=1.
    \end{equation}
    The fit degree is $w_{\mathrm{match}}$, meaningful only beside its
    siblings. Different Concepts are evaluated by separate assessment Events
    and Cuts and need not compete.
    A canonical fragment may itself be opened structurally; component-level
    match judgments use one assessment Event and local Cut per component. No
    independent Realization degree exists. Fit and confidence answer different
    questions. When needed, a separate assessment Cut allocates evidential
    confidence among declared rival interpretations and unresolved evidence.
    Its weights never multiply or replace the fit Cut. A strongly supported
    mismatch and a weakly supported match are therefore distinct. A fit
    remainder means unassigned fit under its own question, not automatically
    lack of evidence; missing evidence is reported separately.
  \item[A11 Profile---Decide.] A candidate envelope names its base commit and
    branch, proposed records and links, replacement map, read dependencies,
    applicable contracts and tolerances, and refinement witness when needed.
    This is transaction metadata, not a seventh world record. A branch has
    one active revision per logical identity. A stale base or changed read
    dependency rejects the candidate; rebasing requires renewed validation.
    Reject leaves the active graph unchanged, although a construction log
    may retain the rejected candidate and explanation. Commit appends one
    immutable cross-surface snapshot and advances its head atomically.
    Revision explicitly replaces named active revisions, preserves their
    history, and revalidates dependents under the profile's declared dependency
    rules. Affected derived views and passages must be regenerated or marked
    stale, never silently treated as current. Opaque operators must declare
    their read dependencies. Branches never merge merely because identifiers
    match.
  \item[A12 Core---Parsimony.] A new record form is admitted only when a validator
    must distinguish two cases that differ in admissible targets,
    composition, authority, inference, or validation, and no existing form
    with typed fields can preserve the distinction.
  \item[A13 Profile---Progressive externalization.] At declared purpose and resolution,
    a construction externalizes a record when it must persist across later
    commits, constrain a later Event, differ across perspective processes, or remain controllable
    and auditable. Other detail remains implicit until an opening condition is
    met. If a governing Cut exists, its unallocated share remains in that
    Cut's remainder; no remainder is invented for an unweighted relation. This
    is a stopping rule, not a completeness claim.
\end{description}

\section{Structural context, perspective projections, and estimates}

\begin{description}
  \item[Q1 Context paths.] Authority-parent edges form an acyclic graph and
    obey A3's unique-context rule. Temporal ancestry is queried separately.
    A record query returns its permitted parents and typed paths to the
    governing root, filtered by the requesting principal's access policy
    before traversal or serialization. A rejected private reference does not
    expose the referenced content or its private ancestry. A Cut or slot
    query returns the complete permitted local Cut---parent Event, question,
    unit, keyed siblings, remainder, optional conditioning address, and
    recomposition rule---or denies that view as a whole. A semantic weight
    is never returned alone. An \texttt{about}, evidence, \texttt{renders},
    or Realization target link is not a permission to traverse its target.
  \item[Q2 Projection.] For an actor $a$, the full actor projection
    $\mathcal M_a(t)$ is derived
    from the belief, appraisal, expectation, and other representational Events
    under $a$'s perspective root that are active at $t$, together with the accepted records
    reachable under Q3. Hypothetical targets remain linked from those inner
    Events rather than being promoted into accepted History. The operative
    model $\mathcal M^{*}_a(t)\subseteq\mathcal M_a(t)$ is the subset recruited
    for one decision or action. These projections represent and organize the
    actor's conduct without assigning a separate perspective field to every record.
  \item[Q3 Access.] $\operatorname{Acc}(a,t)$ contains the exact content
    exposed by observations available to $a$, communications delivered to $a$,
    and $a$'s own eligible earlier records, each available by cutoff $t$.
    The profile records observation scope and delivery time; presence merely
    makes an observation possible and does not reveal another actor's inner
    state or every field of a nearby Event. A communication grants its
    delivered content, not its author's whole model. Later evidence and
    constructor knowledge are excluded. Generation of $a$'s speech, choice, or action
    conditions on an explicitly selected $\mathcal M^{*}_a(t)$ and
    $\operatorname{Acc}(a,t)$ only.
  \item[Q4 Estimate.] Observer $o$'s estimate of a record $\theta$ associated
    with actor $a$ is an Event $e_{\mathrm{est}}$ of kind \textsf{estimate},
    under $o$'s perspective root, with an \texttt{about} link to $\theta$ and evidence
    $\subseteq\operatorname{Acc}(o,t)$. An estimate of categorical or
    externally measured world data may use the corresponding typed Event
    payload, with measurement error where applicable. An estimate of a
    semantic composition is instead a Cut $K_{\mathrm{est}}$ on
    $e_{\mathrm{est}}$, compatible with the target Cut's question, unit,
    vocabulary, and remainder; no interpretive vector is stored as Event
    payload. $\theta$ may itself be an estimate; nesting depth is frozen
    (first construction: $\le 2$).
  \item[Q5 Divergence (derived).] For compatible Cuts $K_{\mathrm{est}}$ and
    $K_\theta$,
    \begin{equation}
      \delta=\tv(K_{\mathrm{est}},K_\theta)=\tfrac12\sum_{v\in V\cup\{\bot\}}
      \bigl|w_{\mathrm{est}}(v)-w_\theta(v)\bigr|;
    \end{equation}
    for comparable discrete fields $\delta=\mathbb{1}[\text{est}\ne\theta]$.
    An incompatible claim in $\mathcal M_a$ and accepted History may be queried
    as misperception; incompatible claims in $\mathcal M_a$ and
    $\mathcal M_b$ may be queried as disagreement. Absence, an unopened region,
    or unequal coverage is ignorance or noncoverage, not automatically error.
    $\delta$ is an external query or validation metric computed
    from Cut weights or categorical equality; it is not stored semantic
    content, an Event measurement, or a Cut.
\end{description}

\section{Surface records (inside the construction graph)}

Understanding Nodes and Document Nodes occupy the shared address space of $\mathcal G_\rho$ but
remain outside the six world-record forms of $\mathcal M_\rho$.

\begin{center}
\footnotesize\ttfamily
\begin{tabular}{@{}l@{}}
UnderstandingNode \{ id, type, text, created at, typed parents, typed links, provenance \}\\
DocumentNode \{ id, role, text, discourse position, optional narrator and viewpoint,\\
\qquad typed parents, typed links, provenance \}
\end{tabular}
\end{center}

\begin{description}
  \item[S1] Understanding links: \texttt{about}, \texttt{supports},
    \texttt{contradicts}, \texttt{refines}, \texttt{answers},
    \texttt{learned-from}, \texttt{supersedes}. Document links:
    \texttt{renders} (story passage to Events), \texttt{shaped-by} (to Understanding Nodes),
    document-scoped \texttt{contains} and \texttt{next}.
  \item[S2] An Understanding Node is not an Event and carries no simplex or
    story-time interval. Its perspective and clock follow the unambiguous
    authority-parent paths to its named understanding-process root. Its
    governing interval may be derived from the commits that create and
    supersede it.
  \item[S3] A Document Node with role \textsf{story-passage} is not the record
    it renders. Document
    \texttt{contains} is not the Event containment of Section~2.
  \item[S4] If a physical carrier of a stored thought must itself be modeled,
    the note, document, or other representational token is a Thing; its creation, reading, revision,
    and supersession are Events. Its text is not an Event.
  \item[S5] The Understanding Nodes beneath a named understanding-process root
    and their links form that process's Understanding Graph. A node may discuss
    any addressed record without borrowing that record's authority. Actor-inner
    Events remain the separate perspective projection defined by Q2, while
    story-passage Document Nodes retain what a character or narrator actually
    expresses.
  \item[S6] A Concept's reasoning history is not stored as an opaque Concept
    field. For one understanding root, it is the ordered subgraph of nodes that
    address the Concept identifier, grounding versions, \textsf{define}
    Realizations, evidence, and predecessors. A revision node records the
    alternatives and evidence behind continuity, split, or merge. Different
    roots may preserve incompatible histories.
  \item[S7] The profile chooses cognitive coverage for the whole construction
    or per subgraph. Understanding Nodes may be absent, selective, or dense.
    Absence means that reasoning was not externalized there, not that no
    reasoning occurred. A claimed rationale trace must state which commitments
    require nodes and which links make the trace checkable. Node density never
    changes world authority.
  \item[S8] A rendering route is a profile, not a world record. It specifies a
    source commit and epistemic status, an ordered Event route and open frontier,
    narrator and viewpoint, evidence cutoff, disclosure rule, and target
    story-passage node. Records outside the route's access set remain hidden
    even though they exist elsewhere in $\mathcal G_\rho$.
  \item[S9] Read-back reconstructs the records implied by a passage in a fresh
    context and compares them with the source commit. Any discrepancy is kept
    as an Understanding Node until a candidate revises the passage, revises the
    addressed records, or marks the ambiguity explicitly.
  \item[S10] Decide under A11 applies to a whole cross-surface candidate.
    Acceptance and rejection therefore cannot leave world, rationale, and
    document records at different commit states.
\end{description}

\section{Interface}

\begin{center}
\footnotesize
\begin{tabularx}{\textwidth}{@{}l Y@{}}
\toprule
Operation & Effect on the construction graph $\mathcal G_\rho$ \\
\midrule
$\mathrm{Register}(c\mid x\mid e)$ & stages an addressable Concept, Thing, or Event; under A1 a Thing and its lifecycle enter the same candidate \\
$\mathrm{Bind}(e,r,x,I)$ & stages a Binding subject to A2 \\
$\mathrm{Open}(P)\mid\mathrm{Compose}(\{e_i\})$ & refines a frozen parent, or proposes a new parent over existing Events \\
$\mathrm{Cut}(P,q,\pi,u_K,\mathrm{slots},h)$ & stages keyed answers and remainder subject to A3--A6 and A9 \\
$\mathrm{Decide}(\Delta)\in\{\mathrm{reject},\mathrm{commit},\mathrm{revise}\}$ & A11 \\
$\mathrm{Realize}(\mathrm{purpose},P,\tau,c,g,\mathrm{evidence})$ & stages a Realization subject to A10 \\
\bottomrule
\end{tabularx}
\end{center}
Convenient names (presence, membership, part--whole, update, roll, roster,
render, and read back) specialize these operations or the surface rules and add
no world-record form.

\paragraph{Conservative opening and Close.}
An opening declares a finite witness $W$: the base snapshot, protected coarse
queries and their observation times and resolution, a fine-to-coarse address
map, a typed projection called $\operatorname{Close}_W$, and a comparison
rule for each query. Each protected answer must be reconstructed from the
addressed fine records, not merely read from a cached coarse answer.
Acceptance as refinement requires
\begin{equation}
  \operatorname{Close}_W(\mathcal G_{\mathrm{base}}\oplus\Delta)(q,t,\rho)
  \simeq_{\varepsilon_q}
  \mathcal G_{\mathrm{base}}(q,t,\rho)
  \qquad\text{for every declared }(q,t,\rho).
\end{equation}
Identity, categorical data, and protected links normally require exact
agreement; measured data use same-unit tolerances. Close is a validation
projection within Open and Decide, not a seventh construction operation.
For a Cut, mapped fine shares sum to the addressed coarse share, including
any part opened from the remainder. Conditional child Cuts multiply by the
share at their conditioning address; their sibling sums remain local.
For unweighted structure, Close preserves the named identities, roles, and
coarse relations through a declared projection; it does not add happenings as
though they were weights. For temporal opening it uses A9 and checks residual
feasibility before all periods have been filled. No claim of conservative
refinement is made without such a witness. Passing protects the stated
finite observations, not every possible future query; a failed opening may
instead be proposed as an explicit revision under A11.

The repository's \texttt{examples/refinement-trial} supplies a small
JavaScript reference fixture for these checks. Its purpose is to make the
contract executable and inspectable, not to certify the Rust host or supply
an independent empirical evaluation of narrative quality.

\section{Derived views, permitted patterns, and operators}

Derived views are computed from committed records, carry derivation lineage,
and hold no authority of their own. Permitted patterns constrain how ordinary
records may be combined, while operators transform or propose records; neither
category introduces another record form.

\begin{center}
\footnotesize
\begin{tabularx}{\textwidth}{@{}p{0.22\textwidth}p{0.13\textwidth}Y@{}}
\toprule
Item & Kind & Definition \\
\midrule
Encapsulation, roster, inherited presence & derived view & A7 \\
Actor perspective projection & derived view & Q2 \\
Concept-understanding history & derived view & the Understanding Node subgraph specified by S6, projected over the addressed grounding lineage of Concept $c$ and ordered in its understanding root's clock \\
Divergence & derived diagnostic & Q5 \\
Want satisfaction & derived view & compare a want's desired Event condition or typed world datum with committed state; any graded judgment is a local Cut \\
Capacity and exchange & permitted pattern & an externally measured stock is typed Event data; allocation is a Cut over one declared unit; transfer is a joint Event. Direct reciprocity or conservation requires a shared external unit and accounting rule; subjective energy is not presumed conserved \\
Shock profile & mixed view & unweighted links identify changed processes; typed data record objective changes where available; an optional Cut allocates attributed impact; anticipation comes from a prior direction Cut \\
Life-change allocation & permitted pattern & allocation Cut over one explicitly named attribution unit among affected processes; heterogeneous objective changes remain separately typed and are not standardized into one scalar \\
Problem-solving profile & derived profile & duration-weighted mixture of scene-level module Cuts over framing, deciding, and their opened modes; decision Events retain strategy families used as unweighted Realizations, while recruitment counts are a theory-only cross-check \\
Recognizer $\Gamma_c$ & operator & maps an addressable fragment to an unweighted \textsf{describe} Realization, role alignment, evidence, and optional local match Cuts \\
Instantiator $\iota_c$ & operator & maps a desired concept state and available support to a candidate fragment under Decide \\
\bottomrule
\end{tabularx}
\end{center}

\section{Book-specific person profile (non-normative)}

A person is a Thing whose life is its lifecycle Event. The Book profile opens
that life into state, orientation, problem solving, feeling, and conditional
change using only ordinary records. Its principals share one concern vocabulary
so their Cuts can be compared, and the remainder tests that vocabulary's
coverage. Another construction may use different questions, different answers,
or no person profile. Only the Cut contract is general.

\begin{description}
  \item[P1 State (IS).] The Book tracks body, kin, partnership,
    work, place, means, knowledge, standing, and meaning as concurrent
    Event-processes within a lifecycle. They may overlap and do not sum.
    Shorter Events may occur within, constitute, or update them through
    unweighted links. A causal interpretation belongs in an Understanding Node
    or a separately posed attribution Cut. Money, location, dates, counts, and
    physical measurements may be Event data; no life domain receives a free-
    standing psychological score. World Events and actor-believed Events use
    different ancestry paths.
  \item[P2 Concerns and wants (WHAT).] At scene $s$, a motivational Event
    $e_{s}^{\mathrm{what}}$ contained by the actor's lifecycle parents every Cut in the WHAT tree. Its top
    Cut divides one unit of scene motivational attention among the five
    first-run concerns and remainder. For each opened concern $i$, another Cut
    on the same Event divides the conditional unit ``attention within concern
    $i$'' among approach, threat-avoidance, and remainder:
    \begin{equation}
      \sum_i c_i+r_C=1,\qquad a_i+v_i+r_i=1.
    \end{equation}
    An opened pole uses another Cut on that Event whose unit is ``attention
    within concern $i$ and pole $p$''; it divides among AtomicWant leaves
    $\ell_{ipj}$ and remainder $r_{ip}$ with
    $\sum_j\ell_{ipj}+r_{ip}=1$. The composed path share within
    that scene's declared attention unit is $c_i p_{ip}\ell_{ipj}$, where
    $p_{i,\mathrm{approach}}=a_i$ and
    $p_{i,\mathrm{avoidance}}=v_i$. Top, pole, and leaf remainders remain
    explicit; across opened leaves their combined unresolved path mass is
    $r_C+\sum_i c_i r_i+\sum_{i,p}c_i p_{ip}r_{ip}$. The conditional
    five-concern approach view renormalizes $c_i a_i$ by
    $Z_A=\sum_i c_i a_i$, and the threat-avoidance view renormalizes $c_i v_i$
    by $Z_V=\sum_i c_i v_i$. If its denominator is zero, that conditional view
    is absent or remainder-only rather than fabricated. At most eight leaves
    are live for one person in one scene. Derived conditional views retain
    their conditioning mass so their scope remains explicit.
    Thus no concern or answer label becomes a Cut parent; every local weight is
    interpreted under the conditional unit named on its Event-parented Cut.
    Each opened concern or pole Cut records its conditioning slot address $h$,
    so the path is recoverable without interpreting the label text.
  \item[P3 Appraisal chain.] A threat-avoidance leaf remains a want: its
    endpoint keeps one named threat from obtaining. A different Event records
    belief in that threat, and fear is a fast appraisal that may follow. Fear
    is neither the concern nor the pole. The chain is
    concern $\rightarrow$ threat belief $\rightarrow$ fear $\rightarrow$
    decision or action, with one record per link. Want satisfaction has a
    world-relative view and an actor-believed view.
  \item[P4 Problem solving (HOW).] A scene assessment Event beneath the actor's
    lifecycle parents the HOW Cuts. Its top Cut divides problem-solving attention
    among framing, deciding, and a remainder. Conditional Cuts on that same
    Event divide ``attention within framing'' among concrete, abstract, and a
    remainder, and ``attention within deciding'' among impartial, personal,
    and a remainder. Cross-cut combinations such as abstract and personal may
    coexist. A strategy family is an unweighted \textsf{describe} Realization
    unless a separate Event-parented Cut asks a genuine comparison question.
  \item[P5 Feeling (FEELS).] A scene appraisal Event parents an emotion Cut.
    The Cut is elicited as a check against the actor's concerns, beliefs, and
    Events. In particular, the Book checks whether an active threat belief
    supports fear; it does not calculate nine emotion shares from five concern
    shares.
  \item[P6 Time and conditionality.] Scene profiles are elicited. Episode,
    period, and life profiles are derived only across disjoint intervals under
    their authoritative sequential Cuts. Every result reports coverage and
    leaves uncovered duration unresolved. Situation and participants select
    conditional views. Drift, recovery, consolidation, instability or cascade,
    anticipated transition, and terminal transition label shapes rather than
    impose equations.
  \item[P7 Numbers.] A Cut weight means a local share under its Event, question,
    unit, siblings, and remainder. Event-process links carry no weights. World
    measurements retain their external units on Events. The profile introduces
    no anchored psychological coordinate and no Realization degree; validator
    tolerances and derived diagnostics are not world facts. Compatible slow
    profiles come from duration, not another round of elicitation.
\end{description}

\section{Additional non-normative conventions}

\begin{description}
  \item[Description] An Event may carry a readable account of its own version.
    That text neither changes world state nor becomes an Understanding or
    Document Node. Profiles may require descriptions for material Events while
    leaving unopened or mechanical records terse.
  \item[Solvers inside a person] A later profile could model problem-solving
    modules as Things within a life, with recruitment Events and dispositions.
    The Book does not do this; its HOW Cuts are only the simpler projection.
\end{description}

\section{Book-specific constants (non-normative)}

\begin{description}[style=nextline]
  \item[Concurrent human-life process domains.] body, kin,
    partnership, work, place, means, knowledge, standing, meaning. These are
    first-construction Event-process lenses, not universal categories or
    numerical coordinates.
  \item[Emotion vocabulary $V_{\mathrm{emo}}$.] joy, sadness, fear, anger,
    disgust, surprise, trust, anticipation, neutral.
  \item[First-run concern vocabulary $V_{\mathrm{concern}}$.] belonging,
    competence, autonomy, understanding, well-being. These are shared
    first-run comparison categories with a tested remainder, not universal
    concerns. Each may be opened through approach and threat-avoidance poles.
  \item[Problem-solving modes.] framing (concrete, abstract) and deciding
    (impartial, personal), with a remainder at every opened Cut.
  \item[Event kinds (initial registry).] lifecycle, presence, membership,
    physical-part, perception, belief, appraisal, want, estimate, decision,
    action, interaction, communication, travel, consequence, creation,
    destruction, split, merge, replacement, succession.
  \item[AtomicWant.] A view over Event, Binding, Cut, and provenance
    records:
    \begin{center}
    \footnotesize\ttfamily
    \begin{tabular}{@{}l@{}}
    AtomicWant \{ id, parent motivational Event, actor (derived), interval,\\
    \qquad concern, pole, operation, target,\\
    \qquad target field or Event condition, desired endpoint, horizon, depth,\\
    \qquad leaf cut, local leaf weight,\\
    \qquad evidence, provenance \}
    \end{tabular}
    \end{center}
    The lifecycle ancestry of the parent Event identifies the actor. Operation
    is preserve, create, transform, or end; preserving can include keeping an
    absent outcome absent. A target is a registered Thing or Event and may be
    registered before its lifecycle starts. Concern comes from
    $V_{\mathrm{concern}}$, and pole is approach or threat-avoidance. The
    endpoint is categorical or uses a real world-data unit, never a free-standing
    psychological score. Satisfaction follows by categorical or same-unit
    comparison; uncertain evidence about that comparison is a separate
    assessment, not an overwritten world state. No more than eight leaves are
    opened in one scene. The referenced Cut
    supplies local leaf weight, and multiplying the concern, pole, and leaf
    shares gives the scene path share. Context uses a crossing Cut and conflict
    uses an unweighted link. A care leaf may target another person's process,
    but kindness remains a separate episode-level evaluation under A10.
  \item[Completed Book selections.] The construction used exact integer-tick
    Cut normalization, a recorded $.10$ elicitation grid with supported $.05$
    refinement, flat content-addressed concept bundles, no continuous extents,
    an estimate-depth cap of two, and full scene profiles only for its three
    principals. One concept was opened experimentally; none of these choices is
    a core grammar obligation.
\end{description}

\end{document}
